% C+ frontmatter — characterization-first frame, multi-line evidence stack,
% AUC demoted to L6 consistency check.
% JLEO-Cplus: characterization paper with 6 complementary evidence lines.

\begin{frontmatter}

\title{Loser-Side Concentration in Public Procurement: Characterization, Mechanism, and Six Complementary Lines of Evidence\texorpdfstring{\footnote{This version: May 2026. We thank seminar participants at INSPER and anonymous feedback received during paper preparation. All errors are our own. \textbf{Data and code availability:} replication archive (R, Python, and \LaTeX{} sources) is available upon publication; CADE adjudication records and BEC contract awards are public. \textbf{Funding:} the authors received no specific funding for this study. \textbf{Declarations of interest:} none.}}{}}

\author[insper]{Darcio Genicolo-Martins\corref{cor1}}
\ead{darciogm1@insper.edu.br}
\author[insper]{Paulo Furquim de Azevedo}
\affiliation[insper]{organization={INSPER},
  city={S\~ao Paulo}, country={Brazil}}
\cortext[cor1]{Corresponding author.}

\begin{abstract}
We characterize \emph{loser-side concentration}---a participation-only
construct identifying firms that bid in many tenders without ever
winning---in São Paulo's electronic procurement platform (BEC,
\valYearStart--\valYearEnd, \valBECitems\ tender-items). The
construct is theoretically grounded: under a separating-equilibrium
framework with cover bidders and genuine entrants, the joint
statistic ($\log(1{+}\text{tenders\_count})$, $\mathbf{1}[\text{wins}=0]$)
is a sufficient screen for cover-bidder type given award-record data
(Proposition~\ref{prop:exposure_stmt}, Lemma~\ref{lem:wins_zero_stmt}).
We document six complementary lines of empirical evidence consistent
with the framework: (i) a $\valOLSGeneralPBU$ within-item, within-year, within-PBU conditional price premium
in environments containing $\valFL$ flagged firms; (ii) a
$\valPBUgradientRatio\!\!\times$ gradient in the price premium across
quartiles of procuring-unit size, monotonic and consistent with
$\partial m^*/\partial \theta_k < 0$; (iii) selection of Regime~2
(truncated-normal cover bidding) over Regime~1 (uniform high-dispersion)
under BIC in both modalities, with sharper coordination in pregão
($\hat{\sigma}_c/\hat{\sigma}_g = \valStructPregRatio$ vs.\ $\valStructConvRatio$
in convite); (iv) Bayesian latent-class classification of
bid distributions yields $\Pr(\text{cover} \mid \text{FL}) = \valLCAprobCoverFL$
versus $\valLCAprobCoverNonFL$ for non-FL firms, a $\valLCAratio\!\!\times$
ratio that complements but does not strictly substitute our construct; (v) repeated FL-winner
pairs occur $\valDyadicRatio\!\!\times$ above the permutation
null mean, with the top-10 pairs $\valDyadicTopTenRatio\!\!\times$
above null. As a sixth consistency check, we report that the
construct discriminates within the always-loser stratum against
$\valCobidders$ adjudicated CADE-cobidder labels at AUC $\valAUCFLfirm$;
the construct does not discriminate against the $\valDirectCADE$
direct CADE defendants (AUC $\valAUCdirectStd$) in the broader BEC
firm universe, a structural scope limit. Counterfactual welfare
implications, treated as illustrative bounds rather than causal
estimates, range $\valWelfareCrossfit$--$\valWelfareIV$ across
specifications. The empirical content of the construct is robust
across modal regimes and to elimination of the binary cutoff. We
interpret the construct as a characterization of an institutional
regularity in commodity-heavy procurement; we do not claim general
cartel detection.
\end{abstract}

\end{frontmatter}

\noindent\textbf{Highlights:}
\begin{itemize}\setlength{\itemsep}{1pt}
\item Loser-side concentration: theoretically grounded participation-only construct with separating-equilibrium derivation.
\item Six complementary lines of evidence: price premium, oversight gradient, structural regime, Bayesian latent-class agreement, dyadic permutation, and CADE within-stratum AUC.
\item PBU oversight gradient: $\valPBUgradientRatio\!\!\times$ across quartiles, monotonic, consistent with the framework's market-selection prediction.
\item Bayesian latent-class agreement: $\Pr(\text{cover}\mid\text{FL})=\valLCAprobCoverFL$ vs.\ $\valLCAprobCoverNonFL$ for non-FL ($\valLCAratio\!\!\times$ ratio)—complementary, partial overlap.
\item Counterfactual welfare bounds: \valWelfareCrossfit--\valWelfareIV\ ($\sim$$\valWelfareCFOnePctLow$--$\valWelfareCFOnePctHigh$\ of $\valWelfareDenom$ in BEC spending), illustrative not causal.
\item Construct does not flag $\valDirectCADE$ direct CADE defendants in the BEC universe (AUC $\valAUCdirectStd$); scope is explicit and narrow.
\end{itemize}

\bigskip

\noindent\textbf{Keywords:} bid rigging, cover bidding, public procurement, loser-side concentration, participation-only screen, separating equilibrium

\noindent\textbf{JEL No.:} D44, D73, H57, K21, L41

\bigskip
